% ============================================================
%  00_Initialize.tex – Paket içe aktarımları ve değişken tanımları
% ============================================================

% --- LaTeX Paketleri (OTUMF-BPSablonu standart paket listesi) ---
\usepackage[utf8]{inputenc}
\usepackage[T1]{fontenc}
\usepackage[turkish]{babel}
\usepackage{csquotes}
\usepackage{blindtext}
\usepackage{float}
\usepackage{graphicx}
\usepackage{multirow}
\usepackage{multicol}
\usepackage{rotating}
\usepackage{hyperref}
\usepackage{color}
\usepackage{algorithm}
\usepackage{algorithmic}
\usepackage{enumitem,amssymb}
\usepackage{times}
\usepackage{tikz}
\usepackage[space]{grffile}
\usepackage{array}
\usepackage{pifont}
\usepackage{changepage}
\usepackage{commath}
\usepackage{tabularx}
\usepackage{morefloats}
\usepackage[dvi]{totpages}
\usepackage[tbtags]{mathtools}
\usepackage[square, sort&compress,comma,numbers]{natbib}
\usepackage{ragged2e}
\usepackage[none]{hyphenat}
\usepackage{xr}
\usepackage{afterpage}
\usepackage{xcolor}
\usepackage{listings}
\usepackage{fancyhdr}

\renewcommand{\textfraction}{0.10}
\renewcommand{\topfraction}{0.95}
\renewcommand{\bottomfraction}{0.95}
\renewcommand{\floatpagefraction}{0.35}

\graphicspath{{./Figures/}}
\newlist{todolist}{itemize}{2}
\setlist[todolist]{label=$\square$}
\hypersetup{urlcolor=black, colorlinks=true}
\newcommand{\ie}{\textit{i}.\textit{e}. }
\newcommand{\xmark}{\ding{55}}
\usetikzlibrary{spy,calc,shapes,positioning}
\newcolumntype{L}[1]{>{\raggedright\let\newline\\\arraybackslash\hspace{0pt}}m{#1}}
\newcolumntype{C}[1]{>{\centering\let\newline\\\arraybackslash\hspace{0pt}}m{#1}}
\newcolumntype{R}[1]{>{\raggedleft\let\newline\\\arraybackslash\hspace{0pt}}m{#1}}

\definecolor{dkgreen}{rgb}{0,0.6,0}
\definecolor{gray}{rgb}{0.5,0.5,0.5}
\definecolor{mauve}{rgb}{0.58,0,0.82}

\lstset{frame=tb,
  language=Java,
  aboveskip=3mm,
  belowskip=3mm,
  showstringspaces=false,
  columns=flexible,
  basicstyle={\small\ttfamily},
  numbers=none,
  numberstyle=\tiny\color{gray},
  keywordstyle=\color{blue},
  commentstyle=\color{dkgreen},
  stringstyle=\color{mauve},
  breaklines=true,
  breakatwhitespace=true,
  tabsize=3
}

% ============================================================
%  Proje Değişkenleri
% ============================================================

% --- Başlık ---
\newcommand{\ttitle}{%
  Myora: Yapay Zekâ Destekli Mobil Sağlık ve Beslenme Yönetim Platformu Geliştirilmesi}

% --- Yazar ---
\newcommand{\authorname}{Seyit Ali YORGUN}
\newcommand{\studentnumber}{220205014}

% --- Bölüm / Üniversite ---
\newcommand{\dept}{Yazılım Mühendisliği}
\newcommand{\faculty}{Mühendislik Fakültesi}
\newcommand{\university}{OSTİM Teknik Üniversitesi}
\newcommand{\univcity}{Ankara}

% --- Danışman ---
\newcommand{\supname}{Dr.\,Öğr.\,Üyesi Can GÜLDÜREN}

% --- Tarih ---
\newcommand{\degreeyear}{2026}
\newcommand{\degreemonth}{Şubat}
\newcommand{\degreedate}{\degreemonth~\degreeyear}

% --- Anahtar Kelimeler ---
\newcommand{\keywordsTR}{%
  Yapay Zekâ, Mobil Sağlık Uygulaması, RAG Chatbot, GPT-4o,
  Beslenme Analizi, Kan Tahlili, React, Supabase, PostgreSQL, pgvector}
\newcommand{\keywordsEN}{%
  Artificial Intelligence, Mobile Health Application, RAG Chatbot, GPT-4o,
  Nutrition Analysis, Blood Test, React, Supabase, PostgreSQL, pgvector}
